\subsection*{NT-R05--NT-R07. Frobenius error laws with only Dedekind RH}

Fix a finite Galois extension $L/K$ of number fields, with group $G$.
For a real class function $t$ on $G$, put
\[
 \delta_t=\frac1{|G|}\sum_{g\in G}t(g),\qquad
 \pi_t(x)=\sum_{\substack{N\mathfrak p\le x\\\mathfrak p\ \text{unramified in }L}}
 t(\operatorname{Frob}_{\mathfrak p}),\qquad
 E_t(y)=ye^{-y/2}\bigl(\pi_t(e^y)-\delta_t\operatorname{Li}_2(e^y)\bigr).
\]
All assertions below require only RH for $\zeta_L$; hence $M$ suffices,
and $A$ suffices when $L/\mathbb Q$ is abelian. No assertion is uniform
as the extension varies.

\paragraph{The cited analytic input and its exact field requirement (NT-R05).}
Fiorilli--Jouve, \emph{Distribution of Frobenius elements in families of
Galois extensions}, author arXiv:2001.05428v2, Proposition3.18,
pp.32--34, proves the limiting-distribution theorem by a
Besicovitch $B^2$ approximation. Their Corollary3.17 and
Lemmata3.14--3.15, pp.30--32, supply the cyclic-induction version
which avoids Artin holomorphy. We use its specialization to
$\beta=1/2$, for real $t$. This gives the $B^2$ expansion below;
the full analytic almost-periodicity theorem is a cited classical
input. The precise specialization does not require $L/\mathbb Q$
to be Galois.

Here are all coefficients, with the conjugacy-class normalization
made explicit. For each conjugacy class $C$ choose $g_C\in C$,
put $H_C=\langle g_C\rangle$ and $F_C=L^{H_C}$, and define
\[
 a_t(\rho)=\sum_{C\subset G}\frac{|C|}{|G|}t(C)
       \sum_{\chi\in\widehat H_C}\overline{\chi(g_C)}
             \operatorname{ord}_{s=\rho}L(s,\chi,L/F_C).
 \tag{NT-R05a}
\]
The sum over $C$ means a sum over conjugacy classes. All functions
inside this formula are finite-order Hecke functions, with their
full local factors. The coefficient is independent of the choices
and is equivalently the residue at $\rho$ of the formal logarithmic
derivative attached to $t$.

Indeed, character orthogonality in the cyclic group and induction give
\[
 1_C=\frac{|C|}{|G|}
       \sum_{\chi\in\widehat H_C}\overline{\chi(g_C)}
                \operatorname{Ind}_{H_C}^G\chi.
\]
Artin induction for Euler factors, including inertia invariants,
then proves (NT-R05a). Moreover
$\zeta_L=\prod_{\chi\in\widehat H_C}L(s,\chi,L/F_C)$.
Every factor is holomorphic in the open critical strip. Thus RH for
$\zeta_L$ forces the nontrivial zeros of every factor in this formula
onto the critical line, without invoking holomorphy of any
nonabelian irreducible Artin function. This is the required analytic
input reduction in the cited cyclic argument.

Write $\Gamma$ for the set of distinct positive ordinates of zeros
of $\zeta_L$, and set
\[
 q_t=\frac1{|G|}\sum_{g\in G}t(g^2),\qquad
 b_t=-q_t-2a_t(1/2).
\]
Reality of $t$ gives $a_t(\overline\rho)=\overline{a_t(\rho)}$,
so $b_t$ is real. The cited expansion is
\[
 E_t(y)\ \sim_{B^2}\ b_t-
       2\Re\sum_{\gamma\in\Gamma}
        \frac{a_t(1/2+i\gamma)}{1/2+i\gamma}e^{i\gamma y}.
 \tag{NT-R05b}
\]
The notation means convergence of symmetric height truncations for
the seminorm
\[
 \|f\|_{B^2}^2=\limsup_{Y\to\infty}\frac1Y
                     \int_{\log2}^Y|f(y)|^2\,dy.
\]
The constant $-q_t$ is the prime-square term; the central zero
contributes $-2a_t(1/2)$. The local zero count NT-Z10 implies
\[
 \sum_{T<\gamma\le T+1}|a_t(1/2+i\gamma)|^2
       \ll_{L,K,t}\log^2(T+3),\qquad
 \sum_{\gamma\in\Gamma}
 \frac{|a_t(1/2+i\gamma)|^2}{1/4+\gamma^2}<\infty.
 \tag{NT-R05c}
\]
To check this directly, each nonnegative factor order in
(NT-R05a) is at most the corresponding order of $\zeta_L$.
The sum of these latter orders in a unit interval is $O_L(\log(T+3))$.
Taking absolute values in the finite class-character sum and then
squaring proves the first estimate; summing over integer intervals
proves the second. Multiplicities are retained throughout.

\paragraph{Complete joint-law and covariance deduction (NT-R06).}
For any finite list of real class functions $t_1,\ldots,t_d$, the
vector $(E_{t_1}(y),\ldots,E_{t_d}(y))$ has a limiting probability
law $\nu$ on $\mathbb R^d$. Precisely, the averages of every bounded
continuous function of this vector converge to its $\nu$ integral.
The mean vector is $(b_{t_j})_j$, and its covariance matrix is
\[
 \Sigma_{jk}=2\Re\sum_{\gamma\in\Gamma}
       \frac{a_{t_j}(1/2+i\gamma)
             \overline{a_{t_k}(1/2+i\gamma)}}{1/4+\gamma^2}.
 \tag{NT-R06a}
\]
Every series here converges absolutely. This is a full new
formulation of the joint deduction from the cited scalar expansion,
not a claim of literature novelty.

To prove it, take the closure $\mathcal T$ of the one-parameter
subgroup
$y\mapsto(e^{i\gamma y})_{\gamma\in\Gamma}$ in the countable
product of circles and give it Haar probability measure. Every
continuous character of a finite-dimensional projection has the
form $e^{i\lambda y}$ along this subgroup. Its Cesaro mean is zero
unless $\lambda=0$, in which case it is one. Stone--Weierstrass
therefore proves that the time averages of every continuous
function of finitely many coordinates equal the corresponding
Haar integrals. Distinct positive frequencies are orthogonal in
$L^2(\mathcal T)$; so are a positive frequency and the conjugate of
another positive frequency.

By (NT-R05c) the series
\[
 Z_j=b_{t_j}-2\Re\sum_{\gamma\in\Gamma}
       \frac{a_{t_j}(1/2+i\gamma)}{1/2+i\gamma}z_\gamma
\]
converges in $L^2(\mathcal T)$. Let $\nu$ be the law of this vector.
For any bounded Lipschitz test function, Cauchy--Schwarz bounds the
difference between its time average on $E$ and on a common finite
truncation by the Euclidean sum of the $B^2$ approximation errors.
The same bound holds for its Haar integrals on $Z$ and that
truncation. The finite-truncation averages converge by the preceding
character argument; letting the truncation grow proves convergence
for bounded Lipschitz tests. These tests determine weak convergence
of probability measures on $\mathbb R^d$, which gives the stated
bounded-continuous formulation. Taking Haar means and second moments
in the $L^2$ series proves (NT-R06a), including the factor two.
Cauchy--Schwarz and (NT-R05c) prove absolute convergence of its
mixed series.

No independence of the coordinates $z_\gamma$ was used: all rational
relations between the ordinates remain encoded in $\mathcal T$.
In particular this argument proves neither a Gaussian law nor
independence of zeros. A zero covariance in some direction means
exactly that its noncentral coefficients $a_t(1/2+i\gamma)$ all
vanish; its limiting law is then the point mass at $b_t$.

\paragraph{Finite etale fibres have a negative mean and a quantitative
one-sided lower density (NT-R07).}
Let $X\to U\subset\operatorname{Spec}\mathcal O_K$ be a nonempty
finite etale scheme, with generic algebra $\prod_{i=1}^r K_i$ and
Galois splitting field $L/K$. Put
\[
 \Pi_X(x)=\sum_{\substack{\mathfrak p\in U\\N\mathfrak p\le x}}
                         \#X(\mathcal O_K/\mathfrak p),\qquad
 E_X(y)=ye^{-y/2}\bigl(\Pi_X(e^y)-r\operatorname{Li}_2(e^y)\bigr).
\]
Apply NT-R05 to the permutation character $t$ of the generic
geometric fibre $S$. Omitting the finitely many additional primes
outside $U$ changes $\pi_t(x)$ by a bounded amount; its normalized
contribution tends to zero and has $B^2$ seminorm zero. Burnside's
lemma gives $\delta_t=r$. The full Artin permutation identity gives
\[
 a_t(1/2)=\sum_{i=1}^r\operatorname{ord}_{s=1/2}\zeta_{K_i}(s)
       \ge0.
\]
To evaluate the sign of $q_t$, decompose $S=\coprod_iG/H_i$.
Counting the pairs $(g,aH_i)$ fixed by $g^2$ gives
\[
 q_t=\sum_{i=1}^r
      \frac{\#\{g\in G:g^2\in H_i\}}{|H_i|}\ge r,
\]
because every element of $H_i$ belongs to the set being counted.
Thus the mean is $-B_X$, where
\[
 B_X=q_t+2\sum_i\operatorname{ord}_{1/2}\zeta_{K_i}\ge r>0,
 \qquad
 V_X=2\sum_{\gamma\in\Gamma}
          \frac{|a_t(1/2+i\gamma)|^2}{1/4+\gamma^2}<\infty.
\]
Consequently
\[
 \liminf_{X_0\to\infty}\frac1{\log X_0}
 \int_2^{X_0}
 1_{\{\Pi_X(x)<r\operatorname{Li}_2(x)\}}\,\frac{dx}{x}
       \ge\frac{B_X^2}{B_X^2+V_X}>0.
 \tag{NT-R07a}
\]
Here is the probability step in full. If $Z$ has the limiting law,
then $W=Z+B_X$ has mean zero and variance $V_X$. For $u\ge0$,
Markov's inequality gives
$\Pr(W\ge B_X)\le(V_X+u^2)/(B_X+u)^2$.
For $V_X>0$, choose $u=V_X/B_X$; for $V_X=0$ the variable
$Z=-B_X$ almost surely. In both cases
$\Pr(Z<0)\ge B_X^2/(B_X^2+V_X)$.
Weak convergence and the open set $(-\infty,0)$ imply the same
lower bound for the liminf of the time proportions. The substitution
$x=e^y$ gives (NT-R07a).

This is a positive lower logarithmic density. The existence of the
sign density, its being greater than one half, and densities of all
strict orderings are not asserted. They require additional control
of the limiting measure, notably its mass on the relevant boundary.
Geometrically these statements concern the trace of Frobenius on
finite-monodromy local systems and, for NT-R07, the permutation
realization of $H^0$ of a finite etale arithmetic cover. They use
ordinary Hecke/Dedekind realizations and do not invoke GRH for an
infinite-image geometric representation.
